% ============================================================================
%  Informe de Proyecto --- Climate Cognitive System (CCS)
%  Curso: Computación Cognitiva --- UTEC
%  Compilar:  tectonic informe.tex        (recomendado, sin dependencias)
%         o:  latexmk -pdf informe.tex
%  Único documento del repositorio en español (el resto está en inglés).
% ============================================================================
\documentclass[11pt,a4paper]{article}

% Engine-agnostic font setup: works with pdflatex AND xelatex/lualatex/tectonic.
\usepackage{iftex}
\ifPDFTeX
  \usepackage[utf8]{inputenc}
  \usepackage[T1]{fontenc}
  \usepackage{lmodern}
  \usepackage{textcomp}
\else
  \usepackage{fontspec}   % defaults to Latin Modern Roman under xetex/lualatex/tectonic
\fi
\usepackage[spanish,es-tabla]{babel}
\usepackage[a4paper,margin=2.3cm]{geometry}
\usepackage{microtype}
\usepackage{booktabs}
\usepackage{array}
\usepackage{enumitem}
\usepackage{xcolor}
\usepackage{tikz}
\usetikzlibrary{shapes.geometric,arrows.meta,positioning,fit,backgrounds,calc}
\usepackage[hidelinks]{hyperref}
\usepackage{titlesec}
\usepackage{float}

\definecolor{ccsblue}{HTML}{1F6FEB}
\definecolor{ccsdark}{HTML}{0B2545}
\definecolor{ccsgray}{HTML}{5B6773}
\definecolor{ccssoft}{HTML}{E8F0FE}

\titleformat{\section}{\normalfont\Large\bfseries\color{ccsdark}}{\thesection.}{0.5em}{}
\titleformat{\subsection}{\normalfont\large\bfseries\color{ccsblue}}{\thesubsection}{0.5em}{}

\setlist{itemsep=1pt,topsep=2pt}
\newcolumntype{L}[1]{>{\raggedright\arraybackslash}p{#1}}

% ---------------------------------------------------------------------------
\begin{document}

% ============================ PORTADA =====================================
\begin{titlepage}
\centering
\vspace*{2.2cm}
{\color{ccsblue}\rule{\linewidth}{2pt}}\\[0.6cm]
{\Huge\bfseries\color{ccsdark} Climate Cognitive System\par}
\vspace{0.3cm}
{\Large Gestión climática predictiva para aulas universitarias\par}
{\color{ccsblue}\rule{\linewidth}{2pt}}\\[1.4cm]

{\large\textbf{Curso:} Computación Cognitiva\par}
{\large\textbf{Docente:} Farfán Madariaga, Jaime Moshe\par}
{\large\textbf{Institución:} UTEC --- Barranco, Lima\par}
\vspace{1.2cm}

{\large\textbf{Integrantes del grupo}\par}
\vspace{0.3cm}
\begin{tabular}{l}
\textcolor{ccsgray}{Balcazar Santa Cruz, Kiara Alexandra} \\
\textcolor{ccsgray}{Gonzales Torres, Aaron Naif} \\
\textcolor{ccsgray}{Surco Vergara, Maria Fernanda} \\
\end{tabular}

\vfill
{\large 2 de julio de 2026\par}
\end{titlepage}

% ============================ ÍNDICE (opcional) ============================
\tableofcontents
\newpage

% ============================ 1. PROBLEMA ==================================
\section{Problema identificado}

En los entornos universitarios, los sistemas de aire acondicionado (A/C) de las
aulas operan tradicionalmente en modo \emph{always-on} durante todo el horario
laboral, sin considerar la ocupación real del espacio ni las condiciones
ambientales del momento.

\paragraph{¿Qué problema existe?} El A/C enfría aulas vacías o con baja
asistencia con la misma intensidad que aulas llenas, porque el control se basa
en el horario y no en la ocupación efectiva. Esto genera un consumo energético
sistemáticamente mayor al necesario.

\paragraph{¿A quién afecta?} A la institución (UTEC), que asume el sobrecosto
energético; al personal de mantenimiento, que carece de datos para optimizar; y
al medio ambiente, por la huella de carbono asociada al consumo eléctrico.

\paragraph{¿Por qué es importante resolverlo?} Los sistemas HVAC representan
cerca del 50\,\% del consumo energético de un edificio no
residencial~\cite{perezlombard2008}, la mayor fracción individual del gasto
eléctrico de un campus. Ajustar la demanda a la
ocupación real reduce costos y emisiones sin sacrificar el confort térmico
recomendado (\textasciitilde 20--24\,°C según ASHRAE~55~\cite{ashrae55}).

\paragraph{¿Qué ocurre actualmente si no se soluciona?} Se mantiene el
desperdicio energético, no existe trazabilidad del uso real de los espacios, y
las decisiones de climatización siguen siendo manuales y reactivas en lugar de
anticipatorias.

% ============================ 2. SOLUCIÓN ==================================
\section{Propuesta de solución}

\textbf{Climate Cognitive System (CCS)} es una plataforma IoT de gestión
climática \emph{predictiva} que ajusta dinámicamente el \emph{setpoint} del A/C
combinando ocupación real, condiciones ambientales y un motor cognitivo de
Machine Learning con \emph{fallback} heurístico.

\paragraph{¿Cómo funcionará?} Sensores físicos reportan temperatura, humedad,
CO\textsubscript{2} y CO al backend. El sistema estima la \textbf{ocupación real}
a partir del CO\textsubscript{2} (balance de masa) y la contrasta con la
\textbf{ocupación esperada} de la reserva. Con ambas señales, más el clima
exterior, un modelo predictivo calcula cuántos grados pre-compensar y emite una
acción cognitiva sobre el A/C (\texttt{ON}/\texttt{STANDBY}).

\paragraph{¿Qué valor aporta?} Enfría según quién \emph{está} en el aula, no
según quién fue reservado: un aula medio vacía produce poco
CO\textsubscript{2}, el sistema lo detecta y reduce la demanda, ahorrando
energía. Además calcula el \textbf{ROI energético} frente al sistema tradicional.

\paragraph{¿Qué lo diferencia de una solución tradicional?} El control clásico
es \emph{always-on} por horario. CCS es \emph{feed-forward + feedback}:
anticipa con la reserva y corrige en tiempo real con la ocupación medida,
integrando además un asistente conversacional (IBM Watson) y detección de
emergencias por monóxido de carbono.

% ============================ 3. PERFILES ==================================
\section{Perfiles del sistema}

El sistema implementa control de acceso por roles (RBAC) de tres niveles.

\begin{table}[H]\centering\small
\begin{tabular}{@{}L{2.6cm}L{12.2cm}@{}}
\toprule
\textbf{Perfil} & \textbf{Acciones principales} \\
\midrule
\textbf{Usuario} \newline (\texttt{guest} / \texttt{collaborator}) &
Registrarse e iniciar sesión; consultar el estado de aulas y sensores en tiempo
real; crear reservas; recibir alertas de emergencia (CO); consultar el
asistente conversacional; el \texttt{collaborator} puede además controlar el
sensor de un aula \emph{si tiene una reserva activa}. \\
\addlinespace
\textbf{Administrador} \newline (\texttt{admin}) &
Gestionar usuarios (aprobar, rechazar, reactivar); aprovisionar y configurar
aulas y sensores; supervisar la telemetría recolectada; administrar reservas;
revisar reportes de ROI energético; recargar el modelo de IA. \\
\bottomrule
\end{tabular}
\caption{Perfiles obligatorios y sus funciones.}
\end{table}

Las cuentas \texttt{guest} se auto-aprueban; \texttt{collaborator} y
\texttt{admin} quedan en estado \texttt{pending} hasta aprobación manual del
administrador.

% ============================ 4. APLICACIÓN ================================
\section{Aplicación web}

Se desarrolla una \textbf{aplicación web responsive} (React + Vite +
TypeScript), pensada como panel administrativo y de monitoreo. Módulos
principales:

\begin{itemize}
\item \textbf{Inicio de sesión / Registro} --- autenticación JWT y flujo de
aprobación por rol.
\item \textbf{Dashboard global} --- monitor de telemetría en vivo.
\item \textbf{Aulas y detalle de aula} --- estado del A/C en tiempo real,
historial de lecturas, control físico del sensor.
\item \textbf{Reservas} --- CRUD con validación de horario.
\item \textbf{Infraestructura} --- panel dual Aulas $\leftrightarrow$ Sensores
para aprovisionamiento.
\item \textbf{Reporte de ROI} --- KPIs energéticos y comparación
Tradicional vs.\ Cognitivo.
\item \textbf{Gestión de usuarios} --- aprobación de cuentas (solo admin).
\item \textbf{Chat flotante} --- asistente IBM Watson embebido.
\end{itemize}

Cada módulo respeta el RBAC: el menú se filtra según el rol autenticado, en
espejo con la autorización del backend.

% ============================ 5. HARDWARE ==================================
\section{Hardware a utilizar}

La telemetría se captura con sensores de bajo costo sobre un
\textbf{Arduino UNO} que emite las lecturas por puerto serial a un equipo
puente. En la versión actual se cuenta con un nodo sensor funcional instalado
en el aula A403 --- el script del gateway está versionado en el repositorio
(\texttt{hardware/sistema.py}) y su montaje se evidencia en los videos de
demostración enlazados en §\ref{sec:recursos} ---; el resto del despliegue se
modela mediante un simulador y un dataset sembrado reproducible
(ver §\ref{sec:viab}).

\begin{table}[H]\centering\small
\begin{tabular}{@{}L{2.6cm}L{3.4cm}L{4.2cm}L{3.4cm}@{}}
\toprule
\textbf{Dispositivo} & \textbf{Función} & \textbf{Datos que captura} & \textbf{Frecuencia} \\
\midrule
DHT11 & Ambiente térmico & Temperatura (°C), humedad (\%) & \textasciitilde 2 s \\
MQ-135 & Calidad de aire & CO\textsubscript{2} (ppm, proxy) $\rightarrow$ ocupación real & \textasciitilde 2 s \\
MQ-7 & Seguridad & CO (ADC $\rightarrow$ ppm, curva calibrada) & \textasciitilde 2 s \\
Cámara + MediaPipe & Aforo & Conteo de personas por cruce de línea & tiempo real \\
Emisor IR\textsuperscript{*} & Actuador & Comando de \emph{setpoint} al A/C & bajo demanda \\
Arduino UNO & Nodo sensor & Agrega lecturas y las emite por Serial (JSON) & \textasciitilde 2 s \\
\bottomrule
\end{tabular}
\caption{Hardware, función, datos y frecuencia de envío.}
{\footnotesize\raggedright \textsuperscript{*}Actuador de fase 2: en la
versión inicial la acción cognitiva (\texttt{ON}/\texttt{STANDBY}) se registra
y visualiza en la aplicación; el envío IR real al A/C es funcionalidad
opcional (§\ref{sec:alcance}).\par}
\end{table}

\paragraph{Conexión con la aplicación.} El Arduino UNO emite cada
\textasciitilde 2\,s un JSON por Serial (9600 baud) con las lecturas de los
tres sensores. Un script Python en el equipo puente
(\texttt{hardware/sistema.py}) parsea la trama, convierte la lectura ADC del
MQ-7 a ppm reales mediante la curva del \emph{datasheet}, fusiona el aforo
estimado por la cámara (MediaPipe) y publica \texttt{POST /api/v1/sensors/}
en formato JSON (\texttt{sensor\_id}, \texttt{temperature}, \texttt{humidity},
\texttt{co2\_ppm}, \texttt{co\_ppm}) por HTTPS a la instancia desplegada,
autenticado por \texttt{X-API-Key}. El backend persiste la lectura y responde
la acción cognitiva; el actuador IR recibirá el \emph{setpoint} sugerido. El
umbral de CO $>50$\,ppm dispara alertas de emergencia~\cite{epaco}.

% ============================ 6. IA =======================================
\section{Uso de Inteligencia Artificial}

La IA es el núcleo funcional, no un adorno. Opera en varias capas:

\begin{itemize}
\item \textbf{Predicción de carga térmica (ML).} Un
\texttt{HistGradientBoostingRegressor}~\cite{sklearn} predice cuántos grados
pre-compensar, a partir de temperatura, hora, ocupación esperada, ocupación
real, temperatura exterior y metadata física del aula.
\item \textbf{Estimación de ocupación real.} Balance de masa sobre el
CO\textsubscript{2}~\cite{co2occupancy}: \texttt{personas} $\approx$
(CO\textsubscript{2} $-$ base) / ppm-por-persona.
\item \textbf{Fallback heurístico.} Si no hay modelo entrenado, una regla
determinística mantiene el servicio operativo.
\item \textbf{Detección de anomalías.} Umbrales sobre temperatura, humedad y CO
para emergencias.
\item \textbf{Visión por computadora (aforo).} MediaPipe Face Detection en el
nodo cliente cuenta entradas y salidas por cruce de línea virtual; el aforo se
registra para trazabilidad y validación de la estimación por CO\textsubscript{2}.
\item \textbf{Asistente conversacional (NLP).} IBM Watson Assistant vía
\emph{custom extension} sobre la API.
\item \textbf{Reportes automáticos.} Cálculo de ROI energético (kWh y CO\textsubscript{2} evitados).
\end{itemize}

\paragraph{Datos, resultado y beneficio.} Procesa la telemetría histórica de los
sensores más el contexto de reservas; produce una decisión de setpoint y un
reporte de ahorro; el beneficio es doble: energía/costo para el administrador y
confort estable para el usuario.

% ============================ 7. ARQUITECTURA =============================
\section{Arquitectura inicial de la solución}

\begin{figure}[H]
\centering
\begin{tikzpicture}[
  font=\footnotesize,
  node distance=8mm and 12mm,
  box/.style={draw=ccsdark,fill=ccssoft,rounded corners,minimum height=8mm,
              minimum width=22mm,align=center,thick},
  actor/.style={draw=ccsblue,fill=white,rounded corners,minimum height=8mm,
                minimum width=20mm,align=center,thick},
  ia/.style={draw=ccsblue,fill=ccsblue!12,rounded corners,minimum height=8mm,
             minimum width=22mm,align=center,thick},
  db/.style={draw=ccsdark,fill=white,cylinder,shape border rotate=90,
             aspect=0.25,minimum height=11mm,minimum width=20mm,align=center,thick},
  hw/.style={draw=ccsgray,fill=white,rounded corners,minimum height=8mm,
             minimum width=22mm,align=center,thick},
  arr/.style={-{Stealth[length=2mm]},thick,ccsgray}
]
% actors
\node[actor](user){Usuario};
\node[actor,below=of user](admin){Administrador};
% app
\node[box,right=of $(user)!0.5!(admin)$](app){App Web\\React + Vite};
% backend
\node[box,right=of app](api){Backend / API\\FastAPI};
% databases
\node[db,above right=6mm and 12mm of api](mongo){MongoDB\\(lecturas)};
\node[db,below right=6mm and 12mm of api](pg){PostgreSQL\\(negocio)};
% IA
\node[ia,right=34mm of api](ml){Modelo ML\\scikit-learn};
\node[ia,below=of ml](watson){IBM Watson\\Assistant};
% hardware
\node[hw,below=16mm of app](hw){Arduino UNO + Cámara\\DHT11 · MQ135 · MQ7};

\draw[arr](user)--(app);
\draw[arr](admin)--(app);
\draw[arr]([yshift=1mm]app.east)--([yshift=1mm]api.west);
\draw[arr]([yshift=-1mm]api.west)--([yshift=-1mm]app.east);
\draw[arr](api)--(mongo);
\draw[arr](api)--(pg);
\draw[arr](api)--(ml);
\draw[arr]([xshift=3mm]api.south) -- ++(0,-2.35cm) -| (watson.south);
\draw[arr](hw.east)-|(api.south) node[midway,below,black]{\scriptsize POST /sensors/};
\draw[arr](api.south east)-- ++(0,-6mm) -| (hw.east);
\end{tikzpicture}
\caption{Arquitectura: usuarios y administrador operan la app web, que consume
el backend FastAPI; éste persiste en MongoDB/PostgreSQL, invoca el modelo ML y
actúa como \emph{proxy} de IBM Watson, y recibe la telemetría del hardware IoT.}
\end{figure}

% ============================ 8. TECNOLOGÍAS ==============================
\section{Tecnologías propuestas}

\begin{table}[H]\centering\small
\begin{tabular}{@{}L{3.2cm}L{11.6cm}@{}}
\toprule
\textbf{Capa} & \textbf{Tecnologías} \\
\midrule
Frontend & React 18, Vite, TypeScript, Tailwind CSS, Recharts \\
Backend & Python, FastAPI~\cite{fastapi}, SQLAlchemy, Uvicorn; JWT + Argon2 (\texttt{pwdlib}) \\
Bases de datos & PostgreSQL 15 (negocio), MongoDB 6 (telemetría) \\
IA & scikit-learn (\texttt{HistGradientBoostingRegressor}), IBM Watson Assistant, MediaPipe (visión) \\
Hardware & Arduino UNO, DHT11, MQ-135, MQ-7, cámara + MediaPipe, emisor IR \\
Despliegue & Docker + Docker Compose, Nginx \\
\bottomrule
\end{tabular}
\caption{Stack tecnológico por capa.}
\end{table}

% ============================ 9. ALCANCE ==================================
\section{Alcance inicial del proyecto}\label{sec:alcance}

\begin{minipage}[t]{0.5\linewidth}
\textbf{Funcionalidades obligatorias}
\begin{itemize}
\item Ingesta de telemetría y persistencia.
\item Motor cognitivo (ML + heurístico) con acción sobre el A/C.
\item Estimación de ocupación real vía CO\textsubscript{2}.
\item RBAC de 3 roles + autenticación JWT.
\item App web con dashboards, reservas e infraestructura.
\item Detección y alerta de emergencias por CO.
\item Reporte de ROI energético.
\end{itemize}
\end{minipage}\hfill
\begin{minipage}[t]{0.46\linewidth}
\textbf{Funcionalidades opcionales}
\begin{itemize}
\item Migración a nube gestionada (AWS/GCP).
\item Reentrenamiento automático del modelo.
\item Control real del A/C por emisor IR calibrado.
\item Notificaciones por correo.
\item Exportación de reportes a PDF.
\item Integración de API meteorológica real.
\end{itemize}
\end{minipage}

% ============================ 10. VIABILIDAD ==============================
\section{Viabilidad del proyecto}\label{sec:viab}

\begin{table}[H]\centering\small
\begin{tabular}{@{}L{3.4cm}L{11.4cm}@{}}
\toprule
\textbf{Factor} & \textbf{Evaluación} \\
\midrule
Hardware & Sensores de bajo costo y ampliamente disponibles; un nodo funcional ya operativo. \\
Herramientas de IA & scikit-learn e IBM Watson accesibles sin costo relevante. \\
Conocimientos & Equipo con dominio de Python, React y bases de datos. \\
Tiempo & Backend, frontend y motor cognitivo ya implementados en su primera versión. \\
Riesgo principal & \textbf{Disponibilidad limitada de hardware} (un solo sensor). \\
Mitigación & Simulador de telemetría + \textbf{dataset sembrado reproducible} (\texttt{seed\_data/}), que permite validar todo el pipeline sin múltiples sensores. \\
\bottomrule
\end{tabular}
\caption{Análisis de viabilidad y gestión de riesgos.}
\end{table}

% ============================ 11. IMPACTO =================================
\section{Impacto esperado}

\begin{itemize}
\item \textbf{Ambiental.} Reducción del consumo eléctrico del A/C y, por tanto,
de las emisiones de CO\textsubscript{2} asociadas, al climatizar según ocupación
real en lugar de un horario fijo.
\item \textbf{Económico.} Ahorro directo en la factura energética del campus,
cuantificado por el módulo de ROI.
\item \textbf{Educativo / Institucional.} Trazabilidad del uso real de las aulas,
útil para la planificación de espacios de UTEC.
\item \textbf{Seguridad.} Detección temprana de niveles peligrosos de monóxido
de carbono con alertas en tiempo real.
\item \textbf{Automatización.} Sustituye el control manual y reactivo por
decisiones anticipatorias basadas en datos.
\end{itemize}

% ============================ 12. RECURSOS Y ENLACES ======================
\section{Recursos y enlaces del proyecto}\label{sec:recursos}

\begin{itemize}
\item \textbf{Repositorio de código (GitHub):} \\
\href{https://github.com/Kiarosaurus/Climate-Cognitive-System}{https://github.com/Kiarosaurus/Climate-Cognitive-System}
\item \textbf{Instancia desplegada (producción):} \\
\href{https://kiarosaurus.me}{https://kiarosaurus.me}
\item \textbf{Videos de demostración (Software y Hardware) --- Google Drive:} \\
\href{https://drive.google.com/drive/folders/1n6j4DU4aDTeXRAdlSHInQun2IR3O-uuh?usp=sharing}{https://drive.google.com/drive/folders/1n6j4DU4aDTeXRAdlSHInQun2IR3O-uuh}
\end{itemize}

% ============================ REFERENCIAS =================================
\bibliographystyle{plain}
\bibliography{refs}

\vspace{0.6cm}
\noindent\rule{\linewidth}{0.4pt}\\
{\footnotesize La documentación técnica extensa (arquitectura interna, decisiones
de diseño y detalle del pipeline de ML) se encuentra en el anexo
\texttt{TECHNICAL\_DOCUMENTATION.md} del repositorio.}

\end{document}
